% =============================================================================
% EINLEITUNG ZU BAND 6: HOLOGRAFIE, HLV-BRÜCKE UND DIE STRUKTUR DER ZEIT
% (Bd. 6 von 8; Fortsetzung von Bd. 5)
% =============================================================================

\section*{Sechster Band der Dokumentensammlung}

Mit diesem Band wird die Sammlung der Einzeldokumente zur T0-Theorie / FFGFT
über die ursprünglich fünf Bände hinaus fortgeführt. Seit dem Abschluss von
Band~5 (Dok.~262, Mai 2026) ist der Korpus um mehr als 120 Dokumente
gewachsen. Sie werden in den Bänden 6, 7 und 8 in numerischer Reihenfolge
aufgenommen; die Grenze jedes Bandes ist so gelegt, dass ein inhaltlicher
Schnitt mit dem Umfangsmaß der Druckausgaben zusammenfällt.

\subsection*{Inhalt und Schwerpunkt}

Band~6 umfasst die Dokumente 263 bis 310 aus den Monaten Juni und Juli 2026.
Den roten Faden bildet die Frage, wie die kompakte Geometrie $T^4$ in die
beobachtete dreidimensionale, zeitlich offene Welt übergeht und was dabei an
Information, Phase und Zeitstruktur erhalten bleibt. Am Anfang stehen das
holografische Prinzip in fraktaler Lesart, die kosmologische Entartung
zwischen $\Lambda$CDM und FFGFT sowie die Herleitung der CMB-Peak-Verhältnisse
aus dem $T^4$. Es folgt die Projektionskette vom $T^4$ zum Beobachtungsraum,
die über eine Reihe von Vergleichs- und Brückendokumenten zum
Helix--Light--Vortex-Ansatz (HLV) bis zur gerechneten 3-gegen-6-Brücke in
Dok.~285 führt.

Der dritte Abschnitt behandelt den dynamischen Sektor: kinetische Energie als
$\xi$-fraktalen Anteil des statischen Verhältnisses $\tilde T\cdot m=1$, die
Trennung von Betrag und Phase in den vier Fermion-Sektoren und die doppelte
Herkunft des Winkels $\theta=2/9$, dynamisch (Dok.~291) und ikosaedrisch
(Dok.~293). Die Dokumente 295 bis 307 widmen sich der Zeit und der
Information: die Zeit-Windung als logarithmische Spirale, der Umkehrtest, der
Information als Ausgang und nicht als Grundeinheit ausweist, die native
Zeit-Energie-Reziprozität und die Verortung der Zeit im Zustandsraum. Den
Abschluss bilden der kosmische Sektor und das Skalenanker-Problem im Vergleich
mit $\Lambda$CDM sowie die Resonanzgeometrie als übergeordnete Referenz.

\subsection*{Gliederung}

Die 44 Dokumente sind in fünf Teile gegliedert: Holografie, kosmologische
Entartung und Formalismus-Korrespondenzen (Dok.~263--269);
Dimensionswechsel, HLV-Brücke und Skalenrekursion (Dok.~270--285);
dynamischer Sektor, Betrag und Phase, $\theta=2/9$ (Dok.~286--294); Zeit und
Information (Dok.~295--307); kosmischer Sektor und Resonanzgeometrie
(Dok.~308--310).

\subsection*{Zum Lesen der Dokumente}

Wie in den früheren Bänden ist jedes Dokument eigenständig und wird im Text
mit seiner Korpusnummer zitiert. Die Statusmarker [K], [B], [S] und
[SETZUNG] kennzeichnen den erreichten Ableitungsstand; spätere Korrekturen
und geschlossene Brücken sind im Korrekturregister Dok.~190 verzeichnet,
das laufend im Repository weitergeführt wird (siehe Band~5, Einleitung). Wo ein späteres Dokument eine Aussage dieses
Bandes präzisiert oder zurücknimmt, gilt die spätere Fassung. Die Dokumente
308 bis 310 sind ursprünglich als mehrkapitelige Abhandlungen gesetzt; im
Buch erscheinen ihre Kapitel als Abschnitte des jeweiligen Dokuments.

\subsection*{Stand}

Die Dokumente dieses Bandes tragen den Stand ihrer jeweiligen Fassung im
Korpus zum Oktober 2026.
